\documentclass[12pt, french, latin.medieval, a4paper]{report}
% package de langues 
\usepackage[french]{babel}
\usepackage[utf8]{inputenc}
\usepackage[T1]{fontenc}

% package pour référence de liens sur pdf 
\usepackage[pdftex]{hyperref} % liens
\usepackage{lastpage} % numéro de page
\usepackage{caption}

% package de formes de document 
\usepackage[left=2.5cm,
              right=2.5cm,
              top=2.5cm,
              bottom=2.5cm
             ]{geometry}
% pour les tableaux 
\usepackage{array,multirow,makecell} 
\newcolumntype{R}[1]{>{\raggedleft\arraybackslash }b{#1}}
\newcolumntype{L}[1]{>{\raggedright\arraybackslash }b{#1}}
\newcolumntype{C}[1]{>{\centering\arraybackslash }b{#1}}

% package pour la page de couverture 
\usepackage{graphicx}
\usepackage[cyr]{aeguill}
\usepackage{fancyheadings}
\usepackage{amsmath}
\usepackage{amsfonts}
\usepackage{amssymb}
\usepackage{geometry}
\usepackage{eso-pic}
\usepackage{tikz}
\usetikzlibrary{calc}


% nouvelle commande pour la forme 
%\fancyhf[OFC]{\thepage / \pageref{LastPage}}
\rfoot{\thepage / \pageref{LastPage}}
%\setcounter{tocdepth}{3}

% pakages de couleurs du document
\usepackage{color}
\usepackage{xcolor}


% package pour la chimie
\usepackage[version=4]{mhchem}


% package pour la visualisation du code
\usepackage{listings}
\definecolor{darkWhite}{rgb}{0.94,0.94,0.94}
 
\lstset{
  aboveskip=3mm,
  belowskip=-2mm,
  backgroundcolor=\color{darkWhite},
  basicstyle=\footnotesize,
  breakatwhitespace=false,
  breaklines=true,
  captionpos=b,
  commentstyle=\color{red},
  deletekeywords={...},
  escapeinside={\%*}{*)},
  extendedchars=true,
  framexleftmargin=16pt,
  framextopmargin=3pt,
  framexbottommargin=6pt,
  frame=tb,
  keepspaces=true,
  keywordstyle=\color{blue},
  language=python,
  literate=
  {²}{{\textsuperscript{2}}}1
  {⁴}{{\textsuperscript{4}}}1
  {⁶}{{\textsuperscript{6}}}1
  {⁸}{{\textsuperscript{8}}}1
  {€}{{\euro{}}}1
  {é}{{\'e}}1
  {è}{{\`{e}}}1
  {ê}{{\^{e}}}1
  {ë}{{\¨{e}}}1
  {É}{{\'{E}}}1
  {Ê}{{\^{E}}}1
  {û}{{\^{u}}}1
  {ù}{{\`{u}}}1
  {â}{{\^{a}}}1
  {à}{{\`{a}}}1
  {á}{{\'{a}}}1
  {ã}{{\~{a}}}1
  {Á}{{\'{A}}}1
  {Â}{{\^{A}}}1
  {Ã}{{\~{A}}}1
  {ç}{{\c{c}}}1
  {Ç}{{\c{C}}}1
  {õ}{{\~{o}}}1
  {ó}{{\'{o}}}1
  {ô}{{\^{o}}}1
  {Õ}{{\~{O}}}1
  {Ó}{{\'{O}}}1
  {Ô}{{\^{O}}}1
  {î}{{\^{i}}}1
  {Î}{{\^{I}}}1
  {í}{{\'{i}}}1
  {Í}{{\~{Í}}}1,
  morekeywords={*,...},
  numbers=left,
  numbersep=10pt,
  numberstyle=\tiny\color{black},
  rulecolor=\color{black},
  showspaces=false,
  showstringspaces=false,
  showtabs=false,
  stepnumber=1,
  stringstyle=\color{gray},
  tabsize=4,
  title=\lstname,
}

%package pour les algorithmes
\usepackage{algorithm}
\usepackage{algorithmic}
\usepackage{amsmath}

% package pour les tableaux
\usepackage{longtable}
\usepackage{booktabs}

% package pour les figures 
\usepackage{subcaption}

% package de reférencement
% \usepackage[backend=biber,safeinputenc, sorting=none]{biblatex}

% mettre la bibliographie dans la table de matières
%\usepackage[nottoc, notlof, notlot]{tocbibind}
\include{page_couverture.tex}


\author{MBE MBE MINDJANA LOIC HENRI}
\title{GNPS et MDByTE}

\begin{document}
    
    \pagecouverture
    {Outils bio-informatiques GNPS et MADByTE}
    {}
    \tableofcontents
    \vspace{20cm}

    \section{Informations sur la séance}
    \begin{itemize}
        \item Date : \textbf{4 Avril 2023}, \textbf{12h30}(date locale)
        \item Moyen : visioconférence
    \end{itemize}
    
    \section{Sujet de la séance}
    Il était question pendant cette séance que le Proffésseur Waffo, nous présente comment un biologiste voie les données $MS^{2}$, HSQC et TOCSY de manière globale.
    
    \section{Taches réalisées lors de la séance et résumé}
    Lors de la séance nous avons eu par le biais de la présentation faite par le Proffésseur Waffo, une meilleure appréhension globale de la donnée à manipuler.
    Il a été présenté de manière globale quelles sont les méthodes générales pour réaliser de la déréplication de molécules du vivant, dont on peut citer, la spectroscopie par résonnance magnétique, la spectrométrie de masse, la spectroscopie par rayons UV. 

    
    La spectrométrie par résonnance magnétique s'appuie sur le fait que les atomes prépondérants d'une molécule du vivant possèdent chacune un champ magétique (sous une forme précise) dont les plus importantes sont le carbonne $^{13}C$ et le proton $^{13}H$ ou encore hydrogène, ces dernières peuvent avoir des déplacéments chimiques précises après plongement de la molécule dans un champ magnétique, ces déplacements chimiques constituent donc une trace contenant des informations liées à la structure de la molécules. Par contre la spectrométrie de masse s'appuie sur le fait que lorsqu'on ionise une molécule (positivement ou négativement), on obtient des fragments ionisées de la molécule (radicaux) qui possèdent des valeurs de masse sur charge précises qui séront recoltés, ces valeurs constituent également une trace de la molécule de départ qu'on pourrait qualifié d'empreinte globale. Des améliorations des méthodes permettant d'avoir des empreintes plus précises, ont vue le jour comme la spectrométrie $MS^{2}$ pour le spectrométrie de masse, et la spectroscopie HSQC et TOCSY pour la spectroscopie par résonnance magnétique.


    Lorsqu'on parle de données spectrales $MS^{2}$, on fait référence à l'ensemble des données $MS$ des radicaux obtenues après spectrométrie de masse; pour les données TOCSY, on fait référence au croisement(corrélation) des déplaceménts chimiques des protons entre eux pour obtenir des séquencage de protons d'environnement chimique semblable et pour les données HSQC au croisement des déplacéménts chimiques des protons et de atomes de carbonne pour détecter les liaisons directes entre carbonne et proton.

    Ces techniques décritent plus haut permettent d'avoir par utilisation de molécules "référenciées" un réseau moléculaire qui permet, d'une molécule à l'entrée d'avoir des indications structurales sur elle, des indications qui sont obtenus des molécules référentiées qui sont plus proche de lui sur le réseau.
    
    \section{Taches pour la séance prochaine}
    A la fin de la séance Proffésseur Waffo nous a signalé qu'une reunion par visioconférence sera probablement faite pour qu'il nous présente la vue globale d'un biologiste sur le sujet des rée=seaux moléculaires.

\end{document}